\problemname{Beetle}

En matematiskt sinnad skalbagge befinner sig på en tunn gren, som är misstänkt lik en X-axel. På grenen finns $n$ daggdroppar, var och en innehåller $m$ volymsenheter vatten. Daggdropparna befinner sig på heltalskoordinaterna $x_1, x_2, \dots, x_n$, relativt skalbaggen.

Det är varmt idag. Under en tidsenhet kommer en volymsenhet vatten avdunsta från varje vattendroppe. Skalbaggen är törstig, så törstig att om den når fram till en vattendroppe så dricker den upp den på nolltid. Under en tidsenhet kan skalbaggen krypa en längdenhet, men det som bekymrar skalbaggen är huruvida det lönar sig att krypa efter daggdropparna eller inte.

Din uppgift är att skriva ett program, som givet daggdropparnas koordinater, beräknar den \emph{maximala} mängden vatten skalbaggen kan dricka.

\section*{Indata}
Den första raden innehåller heltalen $n$ och $m$ ($1 \le n \le 300$, $1 \le m \le 10^6$).

De följande $n$ raderna innehåller vardera en av heltalskoordinaterna $x_1, x_2, \dots, x_n$ ($-10{,}000 \le x_i \le 10{,}000$).
Varje $x_i$ är unik.

\section*{Utdata}
Skriv ut ett heltal: den maximala mängd vatten som skalbaggen kan dricka.

\section*{Poängsättning}
Din lösning kommer att testas på en mängd testfallsgrupper.
För att få poäng för en grupp så måste du klara alla testfall i gruppen.

\noindent
\begin{tabular}{| l | l | p{12cm} |}
  \hline
  \textbf{Grupp} & \textbf{Poäng} & \textbf{Gränser} \\ \hline
  $1$    & $5$         & $x_i \geq 0$ \\ \hline
  $2$    & $10$        & $n \leq 20$, $m \leq 100$ \\ \hline
  $3$    & $35$        & $n, m \leq 100$ \\ \hline
  $4$    & $30$        & $n \leq 100$  \\ \hline
  $5$    & $20$        & Inga ytterligare begränsningar.  \\ \hline
\end{tabular}
